\section{Реализация на back-end логиката}

След описването на архитектурата и домейн модела е необходимо да се разгледа конкретната имплементация на back-end логиката. В тази глава фокусът е върху действителното поведение на системата, реализирано чрез контролери, услуги, repository компоненти и конфигурация. Анализът следва структурата на проекта и показва как отделните модули взаимодействат, за да реализират регистрация, управление на фирми и продукти, работа с количка, checkout и история на поръчките.

За разлика от предходните глави, които се концентрират върху общите принципи, настоящата глава стъпва директно върху наличния сорс код. Именно това прави изложението особено ценно: тук архитектурните идеи се виждат в тяхната програмна материализация.

\subsection{Конфигурация на основните инфраструктурни компоненти}

Back-end поведението на приложението зависи от няколко инфраструктурни компонента, дефинирани централизирано. Най-важните сред тях са \texttt{ModelMapper} и \texttt{PasswordEncoder}. Те се регистрират като Spring bean обекти и впоследствие се инжектират в service компонентите.

\begin{lstlisting}[caption={Конфигурация на mapper и password encoder},label={lst:bean-config}]
@Configuration
public class ApplicationBeanConfiguration {

    @Bean
    public ModelMapper modelMapper() {
        return new ModelMapper();
    }

    @Bean
    public PasswordEncoder passwordEncoder() {
        return new Pbkdf2PasswordEncoder();
    }
}
\end{lstlisting}

Листинг \ref{lst:bean-config} показва добър пример за централизирано управление на зависимости. \texttt{ModelMapper} подпомага преобразуването между различни модели, а \texttt{Pbkdf2PasswordEncoder} осигурява по-сигурно съхранение на паролите. Макар тази конфигурация да е кратка, тя има ключово значение за правилното функциониране на регистрацията и профилния модул.

\subsection{Потребителски модул}

Потребителският модул е най-обемен от гледна точка на уеб потоци. Той обхваща регистрация, потвърждение на акаунт, вход, изход, извличане на потребителска информация, редакция на профил и изчисляване на стойността на продуктите в количката.

\subsubsection{Регистрация на потребител}

Регистрационният процес започва в \texttt{UserController}, където POST маршрутът за регистрация приема \texttt{UserRegisterBindingModel}. Ако входните данни са коректни от гледна точка на binding и няма вече съществуващ потребител със същия имейл, контролерът делегира работата към \texttt{UserServiceImpl}. Там се създава нов entity обект от тип \texttt{User}, инициализират се колекциите за компании, количка и история, задават се основните атрибути и паролата се хешира преди запис.

След запис на потребителя системата изпраща електронно писмо с линк за потвърждение. Това е съществена част от логиката, защото не позволява автоматичен вход преди потвърждаване на акаунта. От чисто инженерна гледна точка модулът демонстрира добра последователност: контролерът управлява потока, услугата реализира бизнес логиката, а repository слоят извършва persistence операцията.

\begin{lstlisting}[caption={Опростен вид на регистрационната логика},label={lst:user-register}]
@Override
public Optional<UserRegisterBindingModel> register(
        UserRegisterBindingModel input) {
    if (userRepository.findByEmail(input.getEmail()).isPresent()) {
        return Optional.empty();
    }

    User user = new User();
    user.setCompanies(new ArrayList<>());
    user.setProductsCart(new ArrayList<>());
    user.setHistoryProducts(new ArrayList<>());
    user.setFirstName(input.getFirstName());
    user.setLastName(input.getLastName());
    user.setEmail(input.getEmail());
    user.setConfirmed(false);
    user.setPassword(passwordEncoder.encode(input.getPassword()));
    user = userRepository.saveAndFlush(user);

    SimpleMailMessage message = new SimpleMailMessage();
    message.setTo(input.getEmail());
    message.setSubject("Email Confirmation");
    message.setText(
            "http://localhost:8080/users/"
                    + user.getId()
                    + "/confirm-email");
    emailSender.send(message);

    return Optional.of(input);
}
\end{lstlisting}

Листинг \ref{lst:user-register} илюстрира същината на регистрационния процес. В него ясно се вижда връзката между persistence слоя, сигурността на паролите и e-mail интеграцията. Същевременно могат да се забележат и някои ограничения: URL адресът за потвърждение е твърдо зададен към \texttt{localhost:8080}, а подателят и конфигурацията на пощата зависят от локални настройки.

\subsubsection{Вход в системата}

Входът също е реализиран в \texttt{UserController}. Контролерът подава входния модел към потребителската услуга, която намира потребителя по имейл, проверява дали акаунтът е потвърден и сравнява въведената парола с хеша в базата. При успех имейлът се записва в \texttt{HttpSession}. Този подход е относително прост, но напълно достатъчен за демонстрационен проект.

Важно е да се отбележи, че логиката не използва Spring Security filter chain или security context. Вместо това достъпът до защитени ресурси се контролира ръчно чрез проверки за наличен атрибут в сесията. От методологична гледна точка това прави кода лесен за следване, но въвежда по-голям риск от дублиране на проверки и от пропуски при добавяне на нови маршрути.

\subsubsection{Потвърждение на акаунт и управление на профил}

Потвърждението на акаунта е реализирано чрез маршрут \texttt{/users/\{userId\}/confirm-email}. При извикване на този маршрут потребителската услуга намира съответния потребител и маркира профила като потвърден. След това приложението връща екран за вход.

Профилният модул позволява извличане и обновяване на потребителските данни. Логиката включва две важни проверки: дали текущата сесия е активна и дали старата парола съвпада със съхранения хеш. По този начин системата реализира базов контрол върху промяната на чувствителна информация.

\subsection{Модул за фирми}

Управлението на фирмите е реализирано чрез \texttt{CompanyController} и \texttt{CompanyServiceImpl}. От функционална гледна точка този модул играе ролята на подготвителен етап за продуктовия каталог, защото всяка фирма може да бъде използвана при въвеждане на продукти.

Контролерът за фирми е относително компактен. Той осигурява GET маршрут за визуализация на формата и POST маршрут за запис на фирма. При POST обработката се използва \texttt{AddCompanyBindingModel}, който валидира полета като име, адрес, ДДС номер и регистрационен статус. Service имплементацията намира текущия потребител по имейл от сесията, създава entity обект \texttt{Company}, интерпретира текстовото поле за регистрационен признак и инициализира асоциираната колекция от продукти.

Тук се вижда интересна особеност на дизайна: фирмата не е абстрактна част от глобален каталог, а принадлежи към конкретен потребител. Това подсилва идеята, че приложението е по-близо до персонален търговски workspace, отколкото до многостранна публична маркетплейс система.

\subsection{Модул за продукти и каталог}

Продуктовият модул е ключова част от приложението, защото именно чрез него системата придобива своята каталожна функционалност. Този модул е разделен между \texttt{ProductController}, \texttt{ProductsService} и \texttt{ProductServiceImpl}. Освен това разчита на \texttt{CompanyRepository}, \texttt{ProductRepository}, \texttt{UserRepository} и \texttt{ShoppingCartService}.

\subsubsection{Добавяне на продукт}

Добавянето на продукт започва с форма, която очаква име, цена, количество, категория, фирма и URL към изображение. Контролерът валидира входния модел и при успех извиква услугата за продукти. Там първо се намира фирмата по име, след което се създава нов обект \texttt{Product}. Категорията се преобразува чрез \texttt{CategoryEnum.valueOf()}, а останалите полета се попълват директно от входния модел.

\begin{lstlisting}[caption={Създаване на нов продукт в service слоя},label={lst:add-product-service}]
@Override
public void addProduct(
        ProductAddBindingModel input,
        Object email) throws Throwable {
    Company company = companyRepository.findByName(input.getCompany())
        .orElseThrow(() -> new Throwable("Company not found"));

    Product product = new Product();
    product.setUrl(input.getUrl());
    product.setCategory(CategoryEnum.valueOf(input.getCategory()));
    product.setCompany(company);
    product.setPrice(input.getPrice());
    product.setName(input.getName());
    product.setQuantity(input.getQuantity());

    productRepository.save(product);
}
\end{lstlisting}

Листинг \ref{lst:add-product-service} показва добре една типична service операция: извличане на зависим домейн обект, създаване на нов entity обект, попълване на полета и persistence чрез repository интерфейс. Решението е компактно и четимо, което е предимство за учебна разработка.

\subsubsection{Каталог и филтриране по категории}

Каталогът е реализиран чрез маршрут от вида \texttt{/products/catalog/\{category\}}. В контролера се използва \texttt{switch} конструкция, която делегира към различни service методи според категорията. Тези service методи извличат всички продукти и филтрират резултатите по стойността на \texttt{CategoryEnum}. От архитектурна гледна точка това решение е достатъчно просто и разбираемо. То не е най-ефективният възможен подход, тъй като филтрирането се извършва на ниво Java поток върху колекцията от всички продукти, но за ограничен набор от данни е напълно приемливо.

Каталогът включва и клиентско филтриране по ценови диапазон, което използва REST крайната точка за извличане на максимална цена. Това е интересен пример за комбиниране на сървърно рендиране и допълнителна динамика от JavaScript.

\subsubsection{Преглед на детайли и добавяне в количка}

За всеки продукт системата предоставя детайлен екран, от който потребителят може да заяви определено количество за добавяне в количката. Контролерът извиква метод за проверка на наличността, а при положителен резултат делегира към service логиката за реално добавяне на продукта.

В \texttt{ProductServiceImpl} това става чрез намаляване на наличността на продукта и добавяне или обновяване на запис в количката. Ако продуктът вече е намерен в количката, съществуващият запис се увеличава с новото количество; ако не съществува, създава се нов \texttt{ShoppingCartEntity}. Този модел е функционален, но тук се откроява и важна архитектурна слабост: проверката дали продуктът вече е в количката се извършва глобално върху всички записи в количката, а не само върху количката на текущия потребител. Освен това сравнението се базира на името на продукта, а не на комбинация от потребител и идентификатор на продукт. Това е достатъчно сериозно ограничение и следва да бъде изрично отбелязано.

\subsection{Количка и финализиране на поръчка}

Модулът за количка и checkout е реализиран чрез комбинация между \texttt{ProductController}, \texttt{ShoppingCartServiceImpl}, \texttt{UserServiceImpl} и \texttt{HistoryServiceImpl}. Това е една от най-интересните части на back-end логиката, защото координира няколко различни типа операции.

\subsubsection{Преглед на количката}

При отваряне на количката контролерът добавя към модела списък от записи в количката, брой продукти, обща цена и крайна цена с доставка. Тук обаче отново е налице архитектурна особеност: методът \texttt{findAllShoppingCartEntities()} връща всички записи в количката, без филтриране по потребител. Това означава, че при реална многопотребителска среда визуализацията на количката не е строго изолирана на ниво потребител, което е съществено ограничение на текущата имплементация.

Въпреки това самият поток е ясно структуриран и лесен за проследяване. Потребителят вижда събраните продукти, цените и формата за адрес, а след това може да финализира поръчката чрез отделен POST маршрут.

\subsubsection{Checkout логика и запис в историята}

Checkout процесът е реализиран в \texttt{HistoryServiceImpl} и е маркиран като \texttt{@Transactional}. Това е разумно решение, тъй като операцията променя множество състояния: създава нов исторически запис, обновява или изчиства количката на потребителя, изтрива определени записи и премахва продукти с нулева наличност.

\begin{lstlisting}[caption={Същност на checkout операцията},label={lst:checkout-service}]
@Override
@Transactional
public void secureCheckout(
        String address,
        Object email,
        Double totalPrice) throws Throwable {
    User user = userRepository.findByEmail(email.toString())
        .orElseThrow(() -> new Throwable("User not found"));

    HistoryEntity historyEntity = new HistoryEntity();
    historyEntity.setAddress(address);
    historyEntity.setPrice(totalPrice);
    historyEntity.setDate(LocalDate.now());
    historyEntity.setUser(user);
    historyRepository.save(historyEntity);

    setNullToShoppingCartProducts(user.getId());
    user.setProductsCart(new ArrayList<>());
    user.getHistoryProducts().add(historyEntity);
    userRepository.save(user);

    shoppingCartRepository.deleteAllByUser_Id(user.getId());
    deleteAllProductsWithZeroQuantity();
}
\end{lstlisting}

Листинг \ref{lst:checkout-service} е особено показателен. От една страна, той демонстрира добре как service слой може да координира сложна домейн операция. От друга страна, разкрива и опростения характер на модела. Вместо реална поръчка с отделни позиции се създава обобщен \texttt{HistoryEntity}. Освен това при приключване на покупка продуктите с нулева наличност се изтриват напълно, вместо просто да останат в каталога като изчерпани. Това е валиден компромис за учебен проект, но в реална система би било по-удачно продуктите да останат видими с индикация за липса на наличност.

\subsubsection{Изчистване на историята}

Историята може да бъде изчиствана от потребителя. Реализацията последователно премахва асоциациите от страна на потребителя, занулява връзката към потребителя в историческите записи и изтрива вече осиротелите записи. Тази логика демонстрира осъзнаване на релационните зависимости, макар и реализирано чрез опростен процедурен подход.

\subsection{REST слой}

Освен HTML ориентираното поведение системата включва и \texttt{MaxPriceRestController}, който обслужва маршрут \texttt{/api/products/max-price}. Контролерът извиква \texttt{ProductsService.getProductMaxPrice()} и връща \texttt{MaxPriceViewModel} чрез \texttt{ResponseEntity}. Реализацията е минималистична, но добре подбрана. Тя показва, че същият домейн модел може да бъде използван и за машинно четим интерфейс.

От гледна точка на учебната стойност това е важно, защото демонстрира как в едно и също приложение могат да съществуват както класически MVC маршрути, така и REST крайни точки. В последствие именно тази крайна точка се използва от JavaScript логиката в каталога за конфигуриране на клиентския ценови филтър.

\subsection{Роля на repository интерфейсите в реализацията}

Макар repository интерфейсите да са относително кратки, те играят важна роля в back-end имплементацията. Чрез тях услугите извличат потребители по имейл, фирми по име, продукти по идентификатор, записи от количка по потребител и история на поръчки по потребителски идентификатор. Този стил на работа подчертава предимството на Spring Data JPA: много от ежедневните persistence операции могат да бъдат описани с кратки и четими методи.

Същевременно именно на това ниво се вижда и една от възможностите за бъдеща оптимизация. Част от филтрирането в момента се извършва в service слоя след извличане на големи колекции. При по-натоварена система би било по-ефективно част от тази логика да бъде преместена към по-целенасочени repository заявки.

\subsection{Наблюдавани силни страни на back-end реализацията}

На база направения анализ могат да бъдат формулирани няколко положителни характеристики на back-end слоя:

\begin{itemize}[leftmargin=*]
    \item последователно използване на service слой за по-съществените бизнес операции;
    \item отделяне на HTML контролери и REST контролер;
    \item използване на mapper и binding модели за по-добра структурираност на данните;
    \item прилагане на password hashing и e-mail confirmation механизъм;
    \item използване на транзакционност при операция, която променя множество зависими състояния.
\end{itemize}

Тези качества правят системата убедителна като учебен пример и осигуряват добра основа за бъдещо надграждане.

\subsection{Ограничения и технически рискове в back-end слоя}

Наред със силните страни back-end кодът съдържа и няколко важни ограничения, които трябва да бъдат отчетени честно:

\begin{itemize}[leftmargin=*]
    \item контролът на достъпа е реализиран ръчно чрез сесийни проверки в контролерите, без централизирана security архитектура;
    \item checkout, количка и проверката за вече добавен продукт работят с логика, която не е строго изолирана по потребител;
    \item в част от операциите се разчита на твърдо зададени стойности, като URL за потвърждение и конфигурация за локално изпълнение;
    \item продуктите с нулева наличност се изтриват, вместо да преминават в състояние „изчерпан“;
    \item част от домейн логиката използва по-опростени структури, подходящи за учебен контекст, но недостатъчни за зряла търговска система.
\end{itemize}

Тези ограничения не следва да се разглеждат като провал на проекта. Напротив, те показват къде е естествената граница между функционален демонстрационен прототип и production-ready система.

\subsection{Обобщение}

Back-end реализацията на проекта демонстрира ясно приложени принципи на Spring базирана уеб система: контролери за управление на HTTP потока, service слой за домейн логика, repository интерфейси за persistence, mapper за трансформации и инфраструктурна конфигурация за повторно използваеми компоненти. Реализацията е достатъчно богата, за да покрие пълния цикъл на потребителя в онлайн магазин, и достатъчно прозрачна, за да бъде описана и анализирана в академичен документ.

Същевременно направеният преглед показва и конкретни области за подобрение -- най-вече по отношение на сигурността, изолацията на потребителските данни и по-детайлното моделиране на поръчките. Точно тази комбинация от работеща функционалност и ясно видими точки за развитие придава на проекта неговата инженерна стойност. Следващата глава ще разгледа представящия слой, потребителските маршрути и ролята на интерфейса в общата система.
